\documentclass[12pt, a4paper]{article}
\usepackage{amsmath, amssymb}
\usepackage{geometry}
\usepackage{parskip}
\usepackage{hyperref}
\usepackage{listings}
\usepackage{xcolor}
\usepackage{algorithm}
\usepackage{algpseudocode}
\geometry{margin=2.5cm}

\lstset{
  language=Python,
  basicstyle=\ttfamily\small,
  keywordstyle=\color{blue!70!black}\bfseries,
  commentstyle=\color{green!50!black}\itshape,
  stringstyle=\color{orange!80!black},
  numbers=left,
  numberstyle=\tiny\color{gray},
  numbersep=8pt,
  breaklines=true,
  frame=single,
  rulecolor=\color{gray!40},
  backgroundcolor=\color{gray!5},
  showstringspaces=false,
  tabsize=4
}

\title{\textbf{An Example for Bayesian Statistics}\\
\large in Ecological Forest Modelling}
\author{}
\date{}

\begin{document}

\maketitle

\vspace{-14mm}
Bayesian Statistics treats probability as a before-after confidence rather
than from the classical frequentist perspective.

\section{Bayes' Theorem}

Let $\theta$ denote a model parameter (e.g.\ the growth rate of a tree
species), and let $D$ denote the observed data. Bayes' theorem states

\begin{equation} \label{eq:bayes_prop}
    P(\theta \mid D) = \frac{P(\theta) \; P(D \mid \theta)}{P(D)}, \quad \text{where}
\end{equation}

\begin{itemize}
    \item $P(\theta)$ is the \textbf{prior}, our educated belief about $\theta$
          \textit{before} seeing the data,

    \item $P(D \mid \theta)$ is the \textbf{likelihood}, how probable
          it is to see D if $\theta$ is the true value,

    \item $P(\theta \mid D)$ is the \textbf{posterior}, our updated
          belief about $\theta$ \textit{after} seeing the data as
          the goal of Bayesian inference, and

    \item $P(D)$ is the \textbf{marginal likelihood} (or evidence),
          a normalising constant ensuring the posterior is a valid
          probability distribution. In practice it is often intractable,
          motivating computational methods.
\end{itemize}

In general: \textit{posterior} $\propto$ \textit{likelihood} $\times$
\textit{prior}.

\section{A simple introductory example}

Suppose a disease affects 1\% of the population. A test
returns a positive in 90\% of actually affected patients, but
a false positive in 5\% of healthy patients. A randomly chosen person
tests positive. What is the probability they are actually affected?

\begin{itemize}
    \item $\theta$ = the person is ill, $h$ = the person is healthy, $D$ = the test is positive,
    \item $P(\theta) = 0.01$ as baseline disease rate,
    \item $P(D \mid \theta) = 0.90$ is the likelihood that the test is positive
          given the person is ill,
    \item $P(D)$, that a test is positive, is computed via the law of total probability:
\end{itemize}
\vspace{-1mm}
\begin{equation}
    P(D) = P(D \mid \theta)\,P(\theta)
         + P(D \mid h)\,P(h)
         = 0.90 \times 0.01 + 0.05 \times 0.99 = 0.0585.
\end{equation}

Bayes shows the test is unreliable due to a low prior (disease is rare), as
\vspace{-1mm}
\begin{equation}
    P(\theta \mid D)
        = \frac{0.01 \times 0.90}{0.0585}
        \approx 0.154 = 15.4\% \quad \text{as actual probability}.
\end{equation}






%:----------------------------------------------------
\section{A Forest Modelling Example}

Suppose we seek $\theta$ as the estimation for the mean annual diameter
increment of beech trees in a given region. We have $n$ field measurements $D = \{d_1, \ldots, d_n\}$.

\begin{enumerate}
\item \textbf{Prior.} Assume that previous studies suggest increments are
          typically between 0.2 and 0.8 cm/year. We encode this as a
          Gaussian prior:
          \[
              P(\theta) = \mathcal{N}(\mu_0,\, \sigma_0^2)
              = \frac{1}{\sqrt{2\pi\sigma_0^2}}
                \exp\!\left(-\frac{(\theta - \mu_0)^2}{2\sigma_0^2}\right)
          \]
          with $\mu_0 = 0.5$ and $\sigma_0 = 0.15$, chosen so that
          roughly 95\% of the prior mass falls within
          $\mu_0 \pm 2\sigma_0 = [0.2,\, 0.8]$, consistent with the
          literature. These are \textbf{hyperparameters} fixed using domain knowledge
          before seeing any data. Choosing them is called
          \textbf{prior elicitation}, and is one of the genuinely
          subjective steps in Bayesian inference.

    \item \textbf{Likelihood.} Assuming each measurement is
          independently drawn from a normal distribution with mean
          $\theta$ and known variance $\sigma^2$:
          \[
              P(D \mid \theta) = \prod_{i=1}^{n}
              \frac{1}{\sqrt{2\pi\sigma^2}}
              \exp\!\left(-\frac{(d_i - \theta)^2}{2\sigma^2}\right)
          \]

\item \textbf{Posterior.} We start directly from Bayes' theorem:
\[
    P(\theta \mid D) = \frac{P(D \mid \theta)\; P(\theta)}{P(D)}
\]
$P(D)$ does not depend on $\theta$ so:
\[
    P(\theta \mid D) \;\propto\; P(D \mid \theta)\; P(\theta)
\]
Substituting the likelihood and prior from the steps before:
\[
    P(\theta \mid D)
    \;\propto\;
    \left(
      \prod_{i=1}^{n}
      \frac{1}{\sqrt{2\pi\sigma^2}}
      \exp\!\left(-\frac{(d_i-\theta)^2}{2\sigma^2}\right)
    \right)
    \cdot
    \frac{1}{\sqrt{2\pi\sigma_0^2}}
    \exp\!\left(-\frac{(\theta-\mu_0)^2}{2\sigma_0^2}\right)
\]
The prefactors $\frac{1}{\sqrt{2\pi\sigma^2}}$ and
$\frac{1}{\sqrt{2\pi\sigma_0^2}}$ do not depend on $\theta$, so we drop and merge
\[
    P(\theta \mid D)
    \;\propto\;
    \exp\!\left(
        -\frac{1}{2\sigma^2}\sum_{i=1}^{n}(d_i - \theta)^2
        -\frac{(\theta - \mu_0)^2}{2\sigma_0^2}
    \right)
\]
Taking the logarithm to simplify, keeping proportionality
\[
    \ln P(\theta \mid D)
    \;\propto\;
    -\frac{1}{2\sigma^2}\sum_{i=1}^{n}(d_i - \theta)^2
    -\frac{(\theta - \mu_0)^2}{2\sigma_0^2}
\]
Expanding the squares:
\[
    \sum_{i=1}^{n}(d_i - \theta)^2
    = \sum_{i=1}^{n} d_i^2 - 2\theta\sum_{i=1}^{n} d_i + n\theta^2, \quad
    (\theta - \mu_0)^2 = \theta^2 - 2\mu_0\theta + \mu_0^2
\]
Substituting, dropping terms without $\theta$, sorting:
\[
    \ln P(\theta \mid D)
    \;\propto\;
    -\frac{n\theta^2}{2\sigma^2}
    +\frac{\theta \sum_{i=1}^{n} d_i}{\sigma^2}
    -\frac{\theta^2}{2\sigma_0^2}
    +\frac{\mu_0\,\theta}{\sigma_0^2}
    \;\propto\;
    -\frac{\theta^2}{2}
     \underbrace{\left(\frac{n}{\sigma^2} + \frac{1}{\sigma_0^2}\right)}_{:=\; A}
    +\;\theta\;
     \underbrace{\left(\frac{\sum_{i=1}^n d_i}{\sigma^2}
                      + \frac{\mu_0}{\sigma_0^2}\right)}_{:=\; B}
\]
This is a quadratic in $\theta$. A Gaussian
$\mathcal{N}(\mu, \sigma^2)$ has log-density
$-\frac{(\theta-\mu)^2}{2\sigma^2} \propto -\frac{\theta^2}{2\sigma^2}
+ \frac{\mu\,\theta}{\sigma^2}$,
which is also a quadratic in $\theta$ with exactly this structure,
so the posterior must be Gaussian too. To read off its mean and variance we complete the square.
We factor out $-A/2$ from the quadratic:
\[
    -\frac{A}{2}\theta^2 + B\theta
    = -\frac{A}{2}
      \!\left(\theta^2 - \frac{2B}{A}\theta\right)
\]
We then add and subtract $\left(B/A\right)^2$ inside the bracket
to form a perfect square, since
$\theta^2 - \frac{2B}{A}\theta + \left(\frac{B}{A}\right)^2
= \left(\theta - \frac{B}{A}\right)^2$:
\[
    = -\frac{A}{2}
      \!\left(\theta - \frac{B}{A}\right)^{\!2}
      + \frac{B^2}{2A}
\]
The term $\frac{B^2}{2A}$ is again constant in $\theta$ and vanishes
into $\propto$, leaving directly:
\[
    P(\theta \mid D)
    \;\propto\;
    \exp\!\left(-\frac{\left(\theta - B/A\right)^2}{2/A}\right)
\]
Comparing this with the standard Gaussian kernel
$\exp\!\left(-\frac{(\theta-\mu)^2}{2\sigma^2}\right)$,
we can directly read off mean $\mu_\text{post}= B/A$ and variance $\sigma_\text{post}^2 = 1/A$. Using the definitions of A and B:
\[
    \boxed{
    P(\theta \mid D) = \mathcal{N}(\mu_\text{post},\;
    \sigma_\text{post}^2), \quad
    \sigma_\text{post}^2
        = \left(\frac{n}{\sigma^2}+\frac{1}{\sigma_0^2}\right)^{\!-1},
    \quad
    \mu_\text{post}
        = \sigma_\text{post}^2\!
          \left(\frac{\sum_{i=1}^n d_i}{\sigma^2}
               +\frac{\mu_0}{\sigma_0^2}\right)
    }
\]
The posterior precision $1/\sigma_\text{post}^2$ is the sum of the
data precision and the prior precision. The posterior mean is a
precision-weighted average of the observed mean and $\mu_0$: with
little data the prior dominates; with much data the observations
take over.

\end{enumerate}

\subsection{Numerical Example}
Let there be
$n = 5$ field measurements $D = \{0.52, 0.61, 0.55, 0.48, 0.59\}$,
so $\sum_{i=1}^n d_i = 2.75$, known measurement noise $\sigma^2 = 0.01$,
and prior hyperparameters $\mu_0 = 0.5$, $\sigma_0^2 = 0.0225$.
Computing $A$ and $B$:
\[
    A = \frac{n}{\sigma^2} + \frac{1}{\sigma_0^2}
      = \frac{5}{0.01} + \frac{1}{0.0225}
      \approx 544.4, \quad
    B = \frac{\sum_{i=1}^n d_i}{\sigma^2} + \frac{\mu_0}{\sigma_0^2}
      = \frac{2.75}{0.01} + \frac{0.5}{0.0225}
      \approx 297.2.
\]
The posterior variance and mean follow as
\[
    \sigma_\text{post}^2 = \frac{1}{A} \approx \frac{1}{544.4} \approx 0.00184,
    \qquad
    \mu_\text{post} = \frac{B}{A} \approx \frac{297.2}{544.4} \approx 0.546 \text{ cm/year}.
\]
The observed sample mean is $2.75/5 = 0.55$ and the prior mean is $0.50$.
The posterior mean $\mu_\text{post} \approx 0.546$ lies between the two,
pulled slightly toward the prior; as expected with only five observations.
With more data the posterior would converge toward the sample mean.



%:----------------------------------------------------
\section{Why Bayesian Methods Matter in Ecology}

Ecological systems are inherently uncertain, data is often
noisy, but prior knowledge from field ecology is rich. Bayesian methods
are particularly well-suited because the posterior is a full distribution 
rather than a point estimate.

A 95\% credible interval $[a, b]$ means: given the data and prior,
we are 95\% certain that $\theta \in [a, b]$, a much more
intuitive statement than a frequentist confidence interval, which
only asserts that if the experiment were repeated infinitely many
times, 95\% of the intervals constructed by that procedure would
contain the true value of $\theta$.

Bayesian hierarchical
models (also called mixed-effects or multilevel models) let
parameters vary across sites while pooling information across
them, directly expressed as nested priors:
\[
P(\theta_s \mid \mu, \sigma) \; P(\mu) \; P(\sigma)
\]
where $\theta_s$ is the site-specific parameter, and $\mu$,
$\sigma$ are hyperparameters with their own priors.

The marginal
likelihood $P(D)$ can be used to compare competing models
(e.g.\ different tree growth functions) via the Bayes factor:
\[
    \text{BF}_{12} = \frac{P(D \mid \text{Model 1})}
                        {P(D \mid \text{Model 2})}.
\]

%:----------------------------------------------------
\section{Markov Chain Monte Carlo}

From Bayes' theorem, we had $P(\theta \mid D) = \frac{P(D \mid \theta)\,P(\theta)}{P(D)}$.
For this to be a valid probability distribution, the numerator must
integrate to 1 over all possible values of $\theta$. Enforcing this
normalisation condition gives $P(D)$ explicitly as:
\[
    P(D) = \int P(D \mid \theta) \; P(\theta) \; d\theta
\]
This is simply the law of total probability: we marginalise out
$\theta$ by summing (integrating) the joint probability
$P(D \mid \theta)\,P(\theta)$ over all values it
could take. For a model with many parameters this integral is
intractable: it lives in a high-dimensional space and has no closed
form. Even in our one-parameter example it is only tractable 
because both the prior and likelihood are Gaussian, a
\textbf{conjugacy} that rarely holds in practice.

\textbf{Markov Chain Monte Carlo (MCMC)} sidesteps this problem by only taking
\textit{samples} from it. Suppose we are currently at parameter value $\theta$ and propose
a move to $\theta'$. We want to know whether $\theta'$ is more or
less probable under the posterior. Form the ratio:

\begin{equation}
    r = \frac{P(\theta' \mid D)}{P(\theta \mid D)}
      = \frac{P(D \mid \theta') \; P(\theta')}{P(D \mid \theta) \; P(\theta)}
    \label{eq:ratio}
\end{equation}

$P(D)$ appears in both numerator and denominator and cancels
exactly. We can therefore evaluate $r$ using only the likelihood
and the prior, both of which we can compute directly,
without ever knowing $P(D)$.

\subsection{Monte Carlo Integration}

Many quantities of interest in Bayesian inference are expectations
of the form
\[
    \mathbb{E}[f(\theta)] = \int f(\theta)\, P(\theta \mid D)\, d\theta,
\]
\textbf{Monte Carlo} draws $N$ independent
samples $\theta^{(1)}, \ldots, \theta^{(N)}$ from the target
distribution, replacing the integral with a sample average:
\[
    \mathbb{E}[f(\theta)] \approx \frac{1}{N} \sum_{t=1}^{N} f\!\left(\theta^{(t)}\right).
\]
By the law of large numbers this estimate converges to the true
expectation as $N \to \infty$, regardless of the dimension of
$\theta$. The difficulty is that drawing independent samples from
$P(\theta \mid D)$ directly is also intractable, which is
where Markov chains come in.

\subsection{Markov Chains}

A \textbf{Markov chain} is a sequence of random variables
$\theta^{(1)}, \theta^{(2)}, \ldots$ with the property that each
state depends only on the one immediately before it:
\[
    P\!\left(\theta^{(t+1)} \mid \theta^{(t)}, \theta^{(t-1)}, \ldots\right)
    = P\!\left(\theta^{(t+1)} \mid \theta^{(t)}\right).
\]
Under mild conditions, a Markov chain converges to a unique
\textbf{stationary distribution} $\pi(\theta)$, meaning that once
the chain has run long enough, its samples look as if they were
drawn from $\pi$. The idea behind MCMC is to \textit{construct} a
chain whose stationary distribution is exactly the posterior
$P(\theta \mid D)$, so that running the chain produces samples we
can use for Monte Carlo integration without ever computing
$P(D)$.

\subsection{The Metropolis-Hastings algorithm}

We build a Markov chain whose \textit{stationary distribution} is
the posterior $P(\theta \mid D)$. This is guaranteed by enforcing
\textbf{detailed balance} at every step:
\[
    P(\theta \mid D) \; T(\theta \to \theta')
    = P(\theta' \mid D) \; T(\theta' \to \theta)
\]
where $T$ is the transition kernel. Intuitively, this says that
probability flux from $\theta$ to $\theta'$ must equal the flux
in the reverse direction, so the chain cannot drift away from the
posterior. After a \textit{burn-in} period, during which the
chain travels from its arbitrary starting point towards the high
probability region, subsequent samples are draws from
posterior.

At each step we draw a candidate $\theta'$ from a
\textbf{proposal distribution} $q(\theta' \mid \theta^{(t-1)})$,
which can be any distribution we can sample from. A common choice
is a symmetric Gaussian centred at the current state,
$q(\theta' \mid \theta) = \mathcal{N}(\theta,\, \sigma_\text{prop}^2)$,
so proposals are small random perturbations of the current value.
We then decide whether to accept or reject $\theta'$ by computing
\vspace{-2mm}
\[
    r = \frac{P(D \mid \theta') \; P(\theta')}{P(D \mid \theta^{(t-1)}) \; P(\theta^{(t-1)})}
        \cdot
        \frac{q(\theta^{(t-1)} \mid \theta')}{q(\theta' \mid \theta^{(t-1)})},
\]

where the first factor is the posterior ratio from
equation~\eqref{eq:ratio} and the second factor corrects for any
asymmetry in the proposal: if $q$ is more likely to propose
$\theta'$ from $\theta$ than vice versa, the candidate is
down-weighted to preserve detailed balance.

If $r \geq 1$ the candidate is more probable under the posterior
and is always accepted. If $r < 1$ it is accepted with probability
$r$, allowing the chain to occasionally move to less probable
regions and avoid getting trapped.

\begin{algorithm}[H]
\caption{Metropolis-Hastings}
\begin{algorithmic}[1]
\State Initialise $\theta^{(0)}$ arbitrarily
\For{$t = 1, 2, \ldots, N$}
    \State Propose $\theta' \sim q(\theta' \mid \theta^{(t-1)})$
    \State Compute $r = \dfrac{P(D \mid \theta')\;P(\theta')}
                               {P(D \mid \theta^{(t-1)})\;P(\theta^{(t-1)})}
                        \cdot
                        \dfrac{q(\theta^{(t-1)} \mid \theta')}
                               {q(\theta' \mid \theta^{(t-1)})}$
    \State Draw $u \sim \mathrm{Uniform}(0, 1)$
    \If{$u < \min(1,\, r)$}
        \State $\theta^{(t)} \leftarrow \theta'$ \quad (accept)
    \Else
        \State $\theta^{(t)} \leftarrow \theta^{(t-1)}$ \quad (reject, stay)
    \EndIf
\EndFor
\State \Return $\{\theta^{(1)}, \ldots, \theta^{(N)}\}$
\end{algorithmic}
\end{algorithm}

Because the chain moves step by step, consecutive samples are
\textit{correlated}, so the effective sample size is smaller than
$N$. Two diagnostics help assess whether the chain has
converged and is mixing well:

\begin{itemize}
    \item \textbf{Trace plot}: plot $\theta^{(t)}$ vs.\ $t$. After
          burn-in this should resemble white noise around a stable
          mean, with no visible trend or slow drift.

    \item \textbf{$\hat{R}$ (Gelman-Rubin statistic)}: run multiple
          chains from different starting points and compare their
          variance. $\hat{R} \approx 1$ indicates that all chains have
          converged to the same distribution.
\end{itemize}

%:----------------------------------------------------
\section{MCMC in Practice}

We now apply everything from the previous sections to the beech tree
growth model. The goal is to estimate $\theta$, the mean annual
diameter increment (cm/year), given noisy field measurements and a
prior belief from the literature.

\paragraph{Setup.}
The true mean increment is chosen to be $\theta_\text{true} = 0.52$ cm/year.
We draw $n = 20$ measurements from
$\mathcal{N}(\theta_\text{true},\, \sigma^2)$ with known observation
noise $\sigma = 0.10$. The prior is
$P(\theta) = \mathcal{N}(\mu_0,\, \sigma_0^2)$
with $\mu_0 = 0.40$ and $\sigma_0 = 0.15$.
The prior mean is deliberately offset from the truth to make the
data-update visible, so we will see the posterior pulled away from
$\mu_0 = 0.40$ towards the true value.

We run the Metropolis-Hastings algorithm for $N = 50.000$ iterations
with a symmetric Gaussian proposal
$q(\theta' \mid \theta) = \mathcal{N}(\theta,\, \sigma_\text{prop}^2)$,
$\sigma_\text{prop} = 0.05$. With
$q(\theta' \mid \theta) = q(\theta \mid \theta')$ due to symmetry,
the acceptance ratio is simpler to calculate below.

\paragraph{What to evaluate.}
From Bayes' theorem we have
\[
    P(\theta \mid D) \;\propto\; P(D \mid \theta)\; P(\theta).
\]
 $P(D)$ cancels in the Metropolis-Hastings
acceptance ratio
\[
    r = \frac{P(\theta' \mid D)}{P(\theta \mid D)}
      = \frac{P(D \mid \theta')\; P(\theta')}{P(D \mid \theta)\; P(\theta)}.
\]
At every step we only need to evaluate the \emph{unnormalised}
posterior, i.e.\ the likelihood times the prior, for the proposed
value $\theta'$ and the current value $\theta$.

\paragraph{Log-posterior.}
In practice we work in log-space to avoid floating-point underflow:
with $n = 20$ measurements, the likelihood
$P(D \mid \theta) = \prod_{i=1}^{n} \mathcal{N}(d_i;\, \theta, \sigma^2)$
is a product of 20 small numbers that can underflow to zero in basic
double precision. Taking the logarithm turns the product into a sum.
Dropping all additive constants that do not depend on $\theta$, which
would also cancel in the ratio $r$:
\[
    P(\theta \mid D)
    \;\propto\;
    \underbrace{\exp\!\left(-\frac{1}{2\sigma^2}\sum_{i=1}^{n}(d_i - \theta)^2\right)}_%
    {P(D\mid\theta)\text{: likelihood}}
    \;\cdot\;
    \underbrace{\exp\!\left(-\frac{(\theta - \mu_0)^2}{2\sigma_0^2}\right)}_%
    {P(\theta)\text{: prior}}
\]
\[
    \ln P(\theta \mid D)
    \;\propto\;
    \underbrace{-\frac{1}{2\sigma^2}\sum_{i=1}^{n}(d_i - \theta)^2}_%
    {\ln P(D\mid\theta)\text{: log-likelihood}}
    \;
    \underbrace{-\frac{(\theta - \mu_0)^2}{2\sigma_0^2}}_{\ln P(\theta)\text{: log-prior}}.
\]
The log-likelihood penalises $\theta$ values that are far from the
observed data; the log-prior penalises values far from $\mu_0$.
Their sum is the log of the unnormalised posterior.

\paragraph{Acceptance step.}
At each iteration $t$ the algorithm from Section 5 proposes a candidate
$\theta' \sim \mathcal{N}(\theta^{(t-1)},\, \sigma_\text{prop}^2)$
and decides whether to accept it. Because the proposal is symmetric,
the correction factor $q(\theta^{(t-1)} \mid \theta') /
q(\theta' \mid \theta^{(t-1)})$ equals 1 and drops out.
The log-acceptance ratio therefore simplifies to:
\[
    \ln r = \ln P(\theta' \mid D) - \ln P(\theta^{(t-1)} \mid D).
\]
We then draw $u \sim \mathrm{Uniform}(0,1)$ and accept $\theta'$
whenever $\ln u < \ln r$, which is mathematically identical to
$u < \min(1, r)$ but numerically stable. Intuitively:
\begin{itemize}
    \item If $\theta'$ has a \emph{higher} log-posterior than the current
          state ($\ln r > 0$, so $r > 1$), we always accept and the chain
          moves to a more probable region.
    \item If $\theta'$ has a \emph{lower} log-posterior ($\ln r < 0$,
          so $r < 1$), we still accept with probability $r$. This
          randomness prevents the chain from getting trapped at a local
          mode.
    \item If rejected, the chain stays at $\theta^{(t-1)}$, which is
          also recorded as a sample.
\end{itemize}

\paragraph{Burn-in.}
The first $5\,000$ samples are discarded. Starting from
$\theta^{(0)} = \mu_0 = 0.40$, the chain begins far from the
high-probability region of the posterior. During burn-in, it
\emph{travels} toward that region; these early samples do not
represent the posterior and must be thrown away. After burn-in,
the remaining samples can be treated as draws from
$P(\theta \mid D)$.

\paragraph{Results.}
After discarding burn-in, the $45\,000$ retained samples yield:
\[
    \hat{\mu}_\text{MCMC} \approx 0.517, \qquad
    \hat{\sigma}_\text{MCMC} \approx 0.022, \qquad
    \text{95\% credible interval} \approx [0.474,\; 0.560].
\]
The posterior mean $0.517$ lies between the prior mean $\mu_0 = 0.40$
and the data mean $\approx 0.52$, pulled towards the data because
$n = 20$ observations provide substantial evidence.
The 95\% credible interval is computed from the posterior
$\mathcal{N}(\mu_\text{post},\, \sigma_\text{post}^2)$, since 95\% of
a Gaussian's mass falls within $\pm 1.96$ standard deviations of the mean:
\[
    [\,\mu_\text{post} - 1.96\cdot\sigma_\text{post},\quad
       \mu_\text{post} + 1.96\cdot\sigma_\text{post}\,]
    = [\,0.517 - 1.96\times 0.022,\quad 0.517 + 1.96\times 0.022\,]
    \approx [0.474,\; 0.560].
\]
This means: given the data and prior, there is a 95\% probability
that $\theta$ lies in $[0.474,\, 0.560]$. In the general MCMC case,
where the posterior is not analytically Gaussian, the same interval
is obtained empirically by sorting the retained samples and taking
the 2.5th and 97.5th percentiles. Both approaches agree here because
the posterior happens to be Gaussian.
The acceptance rate is approximately $60\%$, close to the theoretical
optimum of $44\%$ for a one-dimensional Gaussian target.

\paragraph{Analytic verification.}
Because both the prior and the likelihood are Gaussian, the posterior
is also Gaussian - this is the conjugate case derived in Section~3.
It gives us an exact solution to check against:
\[
    \sigma_\text{post}^2
        = \left(\frac{n}{\sigma^2}+\frac{1}{\sigma_0^2}\right)^{\!-1}
        = \left(\frac{20}{0.01}+\frac{1}{0.0225}\right)^{\!-1}
        \approx 4.89 \times 10^{-4},
    \qquad
    \sigma_\text{post} \approx 0.022,
\]
\[
    \mu_\text{post}
        = \sigma_\text{post}^2\!
          \left(\frac{\sum_{i=1}^n d_i}{\sigma^2}
               +\frac{\mu_0}{\sigma_0^2}\right)
        \approx 0.517 \text{ cm/year}.
\]
The MCMC estimates match the analytic values almost exactly,
confirming that the chain has converged. In practice, most ecological
models are \emph{not} conjugate. The likelihood may be non-Gaussian
(e.g.\ count data, spatial dependence), and no
closed-form posterior exists. In those cases MCMC is the only
available approach, and the Metropolis-Hastings algorithm applies
without any modification.

\paragraph{Summary}
MCMC returns a large array of numbers $\{\theta^{(1)}, \ldots, \theta^{(N)}\}$, each a value
of $\theta$ that the chain visited after burn-in. Because the chain is
constructed to visit regions proportionally to their posterior
probability, this array \emph{is} the posterior, represented as a
histogram rather than an equation. Any summary statistic follows by
plain descriptive statistics on that array:
\[
    \hat{\mu}_\text{post} = \frac{1}{N}\sum_{t=1}^{N}\theta^{(t)},
    \qquad
    \hat{\sigma}^2_\text{post}
        = \frac{1}{N}\sum_{t=1}^{N}\!\left(\theta^{(t)}-\hat{\mu}_\text{post}\right)^{\!2},
\]
and the 95\% credible interval is simply the 2.5th and 97.5th
percentiles of the sample array. This is a direct application of
Monte Carlo integration: for any function $f$,
\[
    \mathbb{E}[f(\theta)] \approx \frac{1}{N}\sum_{t=1}^{N}f\!\left(\theta^{(t)}\right).
\]
Setting $f(\theta)=\theta$ gives the mean; setting
$f(\theta)=(\theta-\hat{\mu})^2$ gives the variance. No analytical
derivation is required. This is precisely why MCMC is indispensable
for complex ecological models: even when the posterior has no closed
form, all quantities of interest can be estimated directly from the
samples.

\end{document}